\begin{proof}
\begin{enumerate}
\item \textbf{Vector notation.}  
Define  
\[
\mathbf a=(a_{1},\dots ,a_{n}),\qquad 
\mathbf b=(b_{1},\dots ,b_{n})\in\mathbb R^{n}.
\]  
Their Euclidean inner product and norm are  
\[
\mathbf a\cdot\mathbf b=\sum_{i=1}^{n}a_{i}b_{i},\qquad 
\|\mathbf a\|^{2}= \sum_{i=1}^{n}a_{i}^{2},\qquad 
\|\mathbf b\|^{2}= \sum_{i=1}^{n}b_{i}^{2}.
\]  
Thus we have to show  
\begin{equation}
(\mathbf a\cdot\mathbf b)^{2}\le \|\mathbf a\|^{2}\,\|\mathbf b\|^{2}.
\tag{1}
\end{equation}

\item \textbf{Degenerate case.}  
If $\|\mathbf b\|^{2}=0$ then $b_{i}=0$ for every $i$, so $\mathbf a\cdot\mathbf b=0$ and both sides of \eqref{1} are zero.  
By symmetry the same holds when $\|\mathbf a\|^{2}=0$.  
Henceforth we assume  
\begin{equation}
\|\mathbf a\|^{2}>0,\qquad \|\mathbf b\|^{2}>0.
\tag{2}
\end{equation}

\item \textbf{Quadratic auxiliary function.}  
For $t\in\mathbb R$ set  
\begin{equation}
f(t)=\sum_{i=1}^{n}(a_{i}-t\,b_{i})^{2}.
\tag{3}
\end{equation}  
Each summand is a square, therefore $f(t)\ge0$ for all $t\in\mathbb R$.

Expanding a typical term,
\[
(a_{i}-t b_{i})^{2}=a_{i}^{2}-2t\,a_{i}b_{i}+t^{2}b_{i}^{2},
\]
we obtain  
\[
f(t)=\sum_{i=1}^{n}a_{i}^{2}
      -2t\sum_{i=1}^{n}a_{i}b_{i}
      +t^{2}\sum_{i=1}^{n}b_{i}^{2}.
\]  
Thus  
\begin{equation}
f(t)=A t^{2}+B t+C,
\tag{4}
\end{equation}
with  
\[
A=\sum_{i=1}^{n}b_{i}^{2}>0,\qquad
B=-2\sum_{i=1}^{n}a_{i}b_{i},\qquad
C=\sum_{i=1}^{n}a_{i}^{2}.
\tag{5}
\]

\item \textbf{Discriminant condition.}  
A quadratic $A t^{2}+B t+C$ with $A>0$ is non‑negative for every real $t$ iff its discriminant does not exceed zero:
\begin{equation}
B^{2}-4AC\le0.
\tag{6}
\end{equation}
Applying \eqref{6} to \eqref{4}–\eqref{5} yields the desired inequality.

\item \textbf{Compute the discriminant.}  
Substituting \eqref{5} into \eqref{6} gives  
\begin{align*}
B^{2}-4AC
&=\bigl(-2\sum_{i=1}^{n}a_{i}b_{i}\bigr)^{2}
  -4\bigl(\sum_{i=1}^{n}b_{i}^{2}\bigr)\bigl(\sum_{i=1}^{n}a_{i}^{2}\bigr)\\[2mm]
&=4\Bigl[\Bigl(\sum_{i=1}^{n}a_{i}b_{i}\Bigr)^{2}
          -\Bigl(\sum_{i=1}^{n}a_{i}^{2}\Bigr)
           \Bigl(\sum_{i=1}^{n}b_{i}^{2}\Bigr)\Bigr].
\end{align*}
Dividing the inequality $B^{2}-4AC\le0$ by the positive constant $4$ we obtain exactly \eqref{1}.  
Hence the Cauchy–Schwarz inequality holds.

\item \textbf{Equality condition.}  
Equality in \eqref{1} occurs precisely when the discriminant in \eqref{6} equals zero, i.e.  
\begin{equation}
B^{2}-4AC=0.
\tag{7}
\end{equation}
For a quadratic with $A>0$, discriminant zero means that there exists $t_{0}\in\mathbb R$ such that $f(t_{0})=0$.  
Because $f(t)$ is a sum of non‑negative squares \eqref{3}, $f(t_{0})=0$ forces each term to vanish:
\[
a_{i}-t_{0}b_{i}=0\qquad\text{for every }i,
\]
hence  
\begin{equation}
a_{i}=t_{0}\,b_{i}\qquad(i=1,\dots ,n).
\tag{8}
\end{equation}
Thus $\mathbf a=\lambda\mathbf b$ with $\lambda=t_{0}$.  

Conversely, if $a_{i}=\lambda b_{i}$ for some $\lambda\in\mathbb R$, then
\[
\sum a_{i}b_{i}= \lambda\sum b_{i}^{2},\qquad
\sum a_{i}^{2}= \lambda^{2}\sum b_{i}^{2},
\]
and a direct substitution shows
\[
\bigl(\sum a_{i}b_{i}\bigr)^{2}
= \lambda^{2}\bigl(\sum b_{i}^{2}\bigr)^{2}
= \bigl(\sum a_{i}^{2}\bigr)\bigl(\sum b_{i}^{2}\bigr).
\]
The degenerate case $\lambda=0$ corresponds to both vectors being zero.  
Therefore equality in the Cauchy–Schwarz inequality holds \emph{iff} there exists $\lambda\in\mathbb R$ such that $a_{i}=\lambda b_{i}$ for all $i$.

\end{enumerate}
\end{proof}